% SPDX-License-Identifier: Apache-2.0
% covers: Vreteno
\datasheet{Vreteno}{A three-stage RV32IMC core with machine-mode traps and an AXI host port.}

\dsfacts{\code{cpu/vreteno/src/core.rs}}{\code{vreteno32::core::Vreteno}}%
{\code{//docs:vreteno}, \code{//docs:noc}, \code{//docs:tapeout}}

\subsection*{Function}

\code{Vreteno} is a RISC-V core for the RV32I base set, the M
extension and the C extension's compressed instructions. It is one
process with three stages. The fetch stage reads the instruction at
the program counter into the instruction register, expanding a
compressed one into the thirty-two bit instruction it stands for. The
execute stage decodes that word, computes, and issues
loads and stores on the bus. The writeback stage extends a loaded
word, writes the register file and retires the instruction.

The core holds two memories: the register file, 32 words, and the
instruction memory, \code{IMEM\_WORDS} or 1024 words.
\code{Vreteno::with} loads a program into the instruction memory. The
data memory and every device are on the bus, outside the core.

The bus side is an AXI host written as hardware. The core issues a
burst on \code{issue}, sends a store's beat on \code{wbeat}, and hands
identifiers back on \code{release}. It receives \code{grant},
\code{done} and \code{rdata} from its tracker, an \code{AxiHost}. The
other inputs are \code{rst}, the external interrupt \code{irq}, the
timer's line \code{tirq} and the software interrupt \code{sirq}, which
the controller raises from \code{msip}; and the debug module's, the
halt and resume requests \code{haltreq} and \code{resumereq} and its
register access, \code{dbg\_regno}, \code{dbg\_wdata} and
\code{dbg\_we}, honoured in debug mode (issue 154). The outputs are
\code{halt}, \code{instr}, the word retiring, \code{wb}, a
\code{Writeback} of \code{done}, \code{rd} and \code{val}, \code{dbg},
debug mode, and \code{dbg\_rdata}, the register the module asked for.

\subsection*{Parameters}

\begin{center}\footnotesize
\begin{tabular}{@{}lp{0.7\textwidth}@{}}
\toprule
Parameter & Meaning \\
\midrule
\code{IW} & The width of the AXI identifier on \code{rdata}, \code{done},
\code{grant} and \code{release}, in bits. The tables use 2, as the
board and the runs do. \\
\bottomrule
\end{tabular}
\end{center}

\subsection*{Register map}

The control registers, by CSR number. Software reaches them through
the six CSR instructions. They are not on the bus.

\begin{center}\footnotesize
\begin{tabular}{@{}llp{0.6\textwidth}@{}}
\toprule
Number & Register & What the core keeps \\
\midrule
\code{0x300} & \code{mstatus} & Bits 3 and 7 only; a write is masked
with \code{0x88}. \\
\code{0x304} & \code{mie} & Bits 11 and 7 only, the external and the
timer enable. \\
\code{0x305} & \code{mtvec} & The handler's address; a write clears
bits 1 and 0. \\
\code{0x340} & \code{mscratch} & The whole word. \\
\code{0x341} & \code{mepc} & The trap's address; a write clears
bit 0. \\
\code{0x342} & \code{mcause} & The whole word. \\
\code{0x343} & \code{mtval} & The whole word. \\
\code{0x344} & \code{mip} & Bit 11 is stored: \code{irq} sets it and
a write keeps only that bit. Bit 7 reads as \code{tirq}. \\
\code{0x7c0} & \code{mhalt} & Nothing: it reads as zero. A write of an
odd value stops the core when that instruction retires. \\
\code{0x301} & \code{misa} & Nothing: it reads
\code{0x4000\_1104}, which is \code{MXL} of 1 and the letters
\code{C}, \code{I} and \code{M}. Writable and ignored, as the
specification allows. \\
\code{0xf11} & \code{mvendorid} & Nothing: zero, which means
unimplemented. Read only. \\
\code{0xf12} & \code{marchid} & Nothing: zero. Read only. \\
\code{0xf13} & \code{mimpid} & Nothing: zero. Read only. \\
\code{0xf14} & \code{mhartid} & Nothing: zero, this machine's one
hart. Read only. \\
\bottomrule
\end{tabular}
\end{center}

A CSR instruction that names any other number is an illegal word and
traps.

\dstables{Vreteno}

\subsection*{Behaviour}

\begin{itemize}
\item With \code{rst} high, the fetch counter goes to zero, the
instruction register is emptied, and the load wait and the sequencer
are cleared.
\item The fetch reads the instruction memory at bits 11 to 2 of the
program counter and at the word after. The halfword at the counter is
the upper half of its word when bit 1 is set. If its low two bits are
not both set it is a compressed instruction, expanded in the fetch;
otherwise the halfword after it is the instruction's upper half. The
counter moves on by two or four.
\item \code{ir\_c} says the instruction in execute was compressed.
\code{jal} and \code{jalr} link the address two bytes on for a
compressed one and four for the rest.
\item An instruction in execute that reads the register the writeback
stage is about to write gets that value directly. When that value is a
load's, the instruction waits one cycle and reads the register file
instead.
\item A taken branch, \code{jal}, \code{jalr}, a trap and \code{mret}
redirect the fetch. The word fetched in that cycle is squashed, which
costs one bubble.
\item A load or store whose address has any bit set above bit 11 goes
on the bus as a burst of one beat, size 2, \code{INCR}. A store's beat
has the lanes it covers as its strobe.
\item A load or store waits for room on both \code{issue} and
\code{wbeat}, whatever its address.
\item A store is posted and the core runs on. A load sets
\code{dev\_wait} and sits in writeback, holding the pipeline, until
its answer lands in \code{wb\_dev}.
\item When a write response and a read beat arrive together, the
write's identifier goes back first and the read's waits a cycle. The
core drops what it receives on \code{grant}.
\item \code{ecall} traps with cause 11 and \code{ebreak} with cause 3,
which takes its own address as \code{mtval}. A word the core does not
know traps with cause 2, and \code{mtval} takes the word. A halfword
that is no RV32C instruction expands to itself, which no thirty-two bit
instruction is, so it traps with cause 2 and \code{mtval} takes the
halfword. An unaligned load traps with cause 4 and an unaligned store
with cause 6, and \code{mtval} takes the address that was asked for.
Every other trap writes zero to \code{mtval}.
\item On a trap, \code{mepc} takes the instruction's address, bit 7 of
\code{mstatus} takes bit 3, bit 3 is cleared, and the fetch goes to
\code{mtvec}. \code{mret} sets bit 3 from bit 7, sets bit 7, and goes
to \code{mepc}.
\item The external interrupt is ready when bit 11 is set in both
\code{mie} and \code{mip}. The software interrupt is ready when bit 3
of \code{mie} is set and \code{sirq} is high, and the timer's when bit
7 of \code{mie} is set and \code{tirq} is high. Any of the three is
taken when bit 3 of \code{mstatus} is set, in place of the instruction
in execute. The external interrupt, cause \code{0x8000\_000b}, wins
over the software interrupt, cause \code{0x8000\_0003}, which wins over
the timer's, cause \code{0x8000\_0007}.
\item A multiply or divide starts the sequencer and stalls in execute.
A multiply is one step and a division thirty-two. \code{//docs:vreteno}
states that a multiply takes three cycles and a division thirty-three.
An interrupt cancels the sequencer, and the instruction starts it
again after the handler returns.
\item The sequencer does not start while a device load in writeback is
still waiting for its answer.
\item A write of an odd value to \code{mhalt} stops the fetch and
parks the program counter one word past that instruction, which
retires first. \code{halt} rises as it retires, and stays high.
\item \code{stall}, \code{int\_take} and \code{redirect} are wires kept
as fields, so the trace and the netlist show them.
\end{itemize}

\subsection*{Verification}

\begin{itemize}
\item \code{//cpu/vreteno:lockstep\_test} runs the core with the router, the
data memory, the timer and the serial port against the reference model
in \code{model.rs}. After every cycle the program counter, registers
1 to 31, the eight control registers that hold something and the halt
must agree. At the
halt the data memory, the timer's compare and the serial port's bytes
must agree. Its tests are \code{demo\_program},
\code{a\_multiply\_right\_after\_a\_load},
\code{compressed\_instructions}, \code{an\_unaligned\_access\_traps},
and \code{random\_programs}, which
runs 64 seeds of 200 instructions each with both lengths mixed and
counts that the mix is there.
\item \code{//cpu/vreteno:isa\_test} compares the core's expander with
the decoder's on every compressed halfword, and
\code{//cpu/vreteno:compressed\_test} checks the decoder against the
GNU disassembler on all 49\,152 of them.
\item \code{//cpu/vreteno:hello\_test} (\code{rust\_runs\_on\_vreteno})
and \code{//cpu/vreteno:hello\_cc\_test}
(\code{cpp\_runs\_on\_vreteno}) run compiled programs on the machine
of \code{run.rs} and check what the serial line says and that the core
halted.
\item \code{//cpu/vreteno:board\_test} and \code{//soc:demo\_test} run
the core inside \code{Board} and on the lattice.
\item The build simulates \code{Vreteno<2>}'s netlist, with the
demonstration program in its instruction memory, against the trace of
\code{//cpu/vreteno:vreteno} under nvc and Verilator:
\code{//docs:vreteno\_sim\_vreteno\_vreteno\_tb\_test} and
\code{//docs:vreteno\_vsim\_vreteno\_test}.
\item \code{//cpu/vreteno:vreteno\_synth} and
\code{//cpu/vreteno:vreteno\_pnr} synthesise and place the core on an
\code{xc7a200tfbg484-2}; both are manual. \code{//docs:vreteno}
reports 2151 LUTs as logic, 48 LUTs as distributed RAM, 489
flip-flops, two block RAM tiles and four DSP blocks after synthesis. It
reports that the routed core meets a 10~ns clock with 0.296~ns of
setup slack.
\item \code{//cpu/vreteno/asic:vreteno\_cells} maps the core onto the
Nangate45 cells and builds with everything.
\code{//cpu/vreteno/asic:}\allowbreak\code{vreteno\_pnr} places and routes it, and is
manual.
\end{itemize}

\subsection*{Used by}

\code{Board}, as \code{cpu}. The demonstration
\code{//cpu/vreteno:vreteno}, the lockstep test, and \code{run.rs},
which the \code{hello} and \code{hello\_cc} runs call. The system in
\code{soc}, where the core is the host at corner 0,0 of the lattice.

\subsection*{Limits}

A program is at most 1024 words, 4096 bytes. The fetch uses bits 11
to 2 of the program counter, and the word after wraps too, so a
program that runs off the end wraps into itself.
Nothing loads the instruction memory at run time; a netlist gets its
contents through \code{init}.

A load from an address below \code{0x1000} sends no request. The
instruction retires with the word the previous device load left in
\code{wb\_dev}. A store below \code{0x1000} writes nothing.

A trap handler must start at a whole word: \code{mtvec} keeps no
bits 1 and 0. The core has machine mode only and the fourteen control
registers above, in three kinds: eight it keeps, \code{mhalt}, which
is its own, and five that say what the machine is and hold nothing. A
write to one of the four whose number begins \code{0xf} is an illegal
instruction, which is what read only means in the specification;
\code{misa} is writable and ignores what is written. The two counters answer as well:
\code{mcycle} at \code{0xb00} with its high half at \code{0xb80},
and \code{minstret} at \code{0xb02} with \code{0xb82}, each 64 bits
and each writable. \code{mcycle} counts while the core runs and stops
when the halt stops it; \code{minstret} counts what retires. They are
the one part of the core the lockstep test leaves alone, for the
reason \code{//docs:vreteno} gives, and
\code{//cpu/vreteno:counters\_test} checks them instead. It has one load in
flight at a time. On the board, \code{irq} is the interrupt
controller's line.
